% !TeX root = ../example.tex
\section{论文讲解与引用}
\sectiontocpage
\sectioncoverpage[assets/tsinghua_door]{本节演示论文讲解、文献字段与引用的显示规则。}
\subsection{论文讲解}

\begin{paperframe}{论文页面示例}{Zadeh2017TensorFN,Liu2018EfficientLM}[bottom-left][5.5cm]
      \begin{description}[引用]
            \item[正文] 放置论文要点、方法图或结果；每页聚焦一个主要信息。
            \item[页眉] 上行沿用当前小节，下行显示本页标题。
            \item[引用] 文献键关联 \texttt{bibliography/refs.bib}，本页自动在左下角列出两篇文献。
      \end{description}
      \medskip
      下一页展示生成本页的命令；文献仅用于演示引用排版。
\end{paperframe}

\begin{frame}[fragile]{论文页面示例：对应源码}
      \begin{verbatim}
\begin{paperframe}{论文页面示例}
  {Zadeh2017TensorFN,Liu2018EfficientLM}
  [bottom-left][5.5cm]
  % 论文要点、方法图或实验结果
\end{paperframe}
      \end{verbatim}
      \begin{description}[文献键]
            \item[页眉] 自动沿用当前 subsection；论文标题作为本页 frametitle。
            \item[文献键] 逗号分隔多篇文献，自动显示引用。
            \item[位置] 末尾可选 \texttt{top-right}（默认）、\texttt{bottom-left} 或 \texttt{bottom-right}。
            \item[宽度] 位置之后可继续填写引用宽度，例如 \texttt{[5.5cm]} 或 \texttt{[0.35\string\paperwidth]}。
      \end{description}
\end{frame}

\subsection{文献引用}

\begin{frame}{角落引用位置}
      \setcornercitewidth{5.5cm} % 仅本页引用加宽，下一页恢复全局宽度
      \begin{description}[右上]
            \item[右上] 默认位置支持用逗号分隔多篇文献。
            \item[左下] 最后一行引用与页脚上边缘保持设定距离。
            \item[右下] 适合左侧已有正文内容的页面。
      \end{description}

      \cornercite{Zadeh2017TensorFN,Liu2018EfficientLM}
      \cornercite[bottom-left]{Zadeh2017TensorFN}
      \cornercite[bottom-right]{Liu2018EfficientLM}
\end{frame}

\begin{frame}[fragile]{角落引用用法}
      \begin{verbatim}
\setcornercitewidth{5.5cm}
\cornercite{key1,key2}
\cornercite[bottom-left]{key1}
\cornercite[bottom-right]{key2}
      \end{verbatim}
      \begin{description}[引用宽度]
            \item[引用宽度] 调整命令放在引用之前，仅影响当前页。
            \item[文献信息] 在 \texttt{bibliography/refs.bib} 中维护，汇总页自动收集已引用条目。
      \end{description}
\end{frame}

% These rules describe the theme's current biblatex/IEEE configuration.
\begin{frame}{从文献字段到角落引用}
  \begin{description}[\texttt{shorttitle}]
    \item[\texttt{author}] 默认输出作者；姓名缩写和多人省略由引用样式决定。
    \item[\texttt{year}] 默认接在作者后，显示发表年份。
    \item[\texttt{title}] 接在作者、年份后，显示论文完整标题。
    \item[\texttt{shorttitle}] 仅在缺少 \texttt{title} 时作为替代标题。
  \end{description}
  \medskip
  本页右上角使用 Liu 条目，展示未设置 \texttt{usera} 时的默认输出。
  \cornercite{Liu2018EfficientLM}
\end{frame}
\note{区分标准作者、年份和标题字段；没有自定义短格式时，角落引用从标准字段生成。}

\begin{frame}[fragile]{自定义短文本：usera 的作用}
  以下摘自 \texttt{bibliography/refs.bib} 的 Zadeh 条目（其余字段省略）：
  \begin{verbatim}
usera      = {Zadeh et al., EMNLP 2017},
title      = {Tensor Fusion Network for
              Multimodal Sentiment Analysis},
shorttitle = {Tensor Fusion Network},
year       = {2017},
  \end{verbatim}
  \begin{itemize}
    \item 设置 \texttt{usera} 后：角落显示“自定义文本 + 完整标题”。
    \item 此时不会另加作者和年份；年份、会议简称需写在 \texttt{usera} 中。
    \item 本条目有 \texttt{title}，因此不会改用 \texttt{shorttitle}。
  \end{itemize}
  \cornercite{Zadeh2017TensorFN}
\end{frame}
\note{usera 自定义角落引用的短格式，不替换文末数字标签。不要把机构等未核实信息直接填入示例。}

\begin{frame}[fragile]{论文讲解页：补充论文信息}
  在文献条目中补充以下字段，\texttt{paperframe} 会自动追加信息：
  \begin{verbatim}
userb = {<论文署名机构>},
userc = {<会议或期刊，发表年月；预印本填写 arXiv>},
usere = {<引用次数>},
userf = {<统计来源，截至 YYYY-MM-DD>},
  \end{verbatim}
  \begin{itemize}
    \item 尖括号为填写提示；未核实的信息可以省略，对应行自动隐藏。
    \item 信息在同一段中连续显示，以竖线分隔，不额外占用多行。
    \item 普通角落引用与文末参考文献保持原格式。
  \end{itemize}
\end{frame}
\note{机构、发表状态、引用量和查询日期都是可选字段。此处尖括号内容只是格式占位符；引用数必须保留来源和日期。}

\begin{paperframe}[字段实际呈现]{论文信息字段示例}{Zadeh2017TensorFN}[bottom-left][7cm]
  \begin{itemize}
    \item 左下角沿用普通引用的灰色文字、字号和紧凑行距。
    \item 论文标题仍由引用中的 \texttt{title} 自动输出。
    \item 机构、发表 venue 或 arXiv、引用量和统计日期来自 \texttt{userb}、\texttt{userc}、\texttt{usere}、\texttt{userf}。
    \item 本页尖括号内容是待填写占位符，替换后即可用于正式汇报。
  \end{itemize}
\end{paperframe}
\note{查看附加字段的实际效果。此页用于排版演示，隐藏未填字段可以保持正式汇报简洁。}

\begin{frame}{文末参考文献：保留完整出处}
  \begin{description}[\texttt{booktitle}]
    \item[\texttt{author}] 按当前 IEEE 样式显示作者，名字使用首字母。
    \item[\texttt{title}] 显示论文题名；\texttt{usera} 不会替换这部分。
    \item[\texttt{booktitle}] 对 \texttt{@inproceedings} 条目，显示会议论文集名称。
    \item[\texttt{year}] 显示发表年份；页码、DOI 等按条目类型与样式输出。
  \end{description}
  \begin{notebox}[角落与文末的区别]
    角落引用不自动输出会议、页码或 DOI。
    文末页从已引用条目生成完整列表，按首次引用顺序编号；
    \texttt{usera} 不会改变这些数字标签。
  \end{notebox}
\end{frame}
\note{文末参考文献保留标准出处。角落短引用与文末完整条目承担不同作用。}

\begin{frame}[fragile]{文献键与编译流程}
  \begin{verbatim}
\addbibresource{bibliography/refs.bib}  % 导言区加载数据库
\cornercite{Zadeh2017TensorFN}
\referencespage[break][\small][onecolumn]
  \end{verbatim}
  \begin{description}[文献键]
    \item[文献键] \texttt{@inproceedings\{Zadeh2017TensorFN,...\}} 中的
      \texttt{Zadeh2017TensorFN} 是查找标识，不是打印出来的标题。
    \item[作者名] 多位作者在 \texttt{author} 字段内用 \texttt{and} 分隔。
    \item[编译] 使用 XeLaTeX 与 Biber；修改 \texttt{.bib} 后重新运行 \texttt{latexmk}。
  \end{description}
  \medskip
  只存入数据库但未引用的条目，默认不会出现在文末参考文献中。
\end{frame}
\note{按文献资源、正文引用、文末汇总的顺序回顾流程。biblatex 与 biber 由构建入口统一处理。}

